\documentclass[11pt]{article}

\usepackage[margin=1.1in]{geometry}
\usepackage{amsmath,amssymb,amsthm,mathtools}
\usepackage{enumitem}
\usepackage{hyperref}
\usepackage[nameinlink,noabbrev]{cleveref}

\hypersetup{colorlinks=true,linkcolor=blue,citecolor=blue,urlcolor=blue}

\newtheorem{theorem}{Theorem}[section]
\newtheorem{lemma}[theorem]{Lemma}
\newtheorem{corollary}[theorem]{Corollary}
\theoremstyle{definition}
\newtheorem{example}[theorem]{Example}
\theoremstyle{remark}
\newtheorem{remark}[theorem]{Remark}

\DeclareMathOperator{\supp}{supp}
\DeclareMathOperator{\diag}{diag}
\providecommand{\doi}[1]{\href{https://doi.org/#1}{doi:#1}}

\title{The whisker threshold for WARM graphs with equal rates}
\author{Yitzchak Shmalo\\
Einstein Institute of Mathematics, The Hebrew University of Jerusalem\\
Givat Ram, Jerusalem, Israel}
\date{}

\begin{document}
\maketitle

\begin{abstract}
In the WARM model of van der Hofstad, Holmes, Kuznetsov and Ruszel on a finite graph, each vertex fires at Poisson times and reinforces one incident edge, chosen
with probability proportional to the $\alpha$-th power of its count. Holmes and Kleptsyn proved that with equal
firing rates and $\alpha>25$ the infinitely reinforced edges form whiskers, trees of diameter at most three, and
Kriebel lowered the threshold to $17/4$. We prove the same conclusion for every $\alpha>4/3$, and show that for
$4/3<\alpha\le 2$ every component is a star. The constant cannot be lowered, since for $\alpha<4/3$ the triangle keeps
all three edges with positive probability. This answers the whisker problem of Kleptsyn and of Aletti et al., except
at $\alpha=4/3$. For $\alpha\le 2$ the proof evaluates the Hessian of a Lyapunov function on a test vector whose
increments alternate along a path between two edges next to leaves.
\end{abstract}

\noindent\textit{Mathematics Subject Classification (2020):} 60K35, 37C75.\\
\textit{Keywords:} WARM, reinforced urns, graph-based competition, linear stability, whisker.

\section{Introduction}\label{sec:intro}

Let $G=(V,E)$ be a finite simple loopless graph and $\alpha>1$. Each edge $e$ carries a count $N_t(e)$, with
$N_0(e)\in\{1,2,\dots\}$. Each non-isolated vertex carries an independent Poisson clock of rate one, and when the
clock of $v$ rings, an edge $e\ni v$ is chosen with probability
\[
  \frac{N_t(e)^\alpha}{\sum_{f\ni v}N_t(f)^\alpha}
\]
and its count is increased by one. This is the WARM model ($W$, $A$-reinforcement model) with equal firing rates and
power reinforcement. It was introduced, for general
rates and a general reinforcement function, by van der Hofstad, Holmes, Kuznetsov and Ruszel \cite{vdHHKR}, as a
model for the way neurons select among connections. Write
\[
  E_\infty=\{e\in E:\ N_t(e)\to\infty\}.
\]
A \emph{whisker} is a tree of diameter at most three, and a \emph{star} is a tree of diameter at most two (a single
edge included).

Van der Hofstad, Holmes, Kuznetsov and Ruszel \cite[Conjecture~2]{vdHHKR} conjectured that for every graph and
$\alpha$ large enough depending on the graph, the linearly stable equilibria of the associated flow are supported on
whisker forests. They proved that for the triangle the uniform equilibrium is linearly stable if and only if
$\alpha<4/3$ \cite[Proposition~2(v)]{vdHHKR}, and that for $\alpha>4/3$ every linearly stable equilibrium of the
triangle has a zero coordinate \cite[Theorem~4(ii)]{vdHHKR}. Holmes and Kleptsyn \cite{HK} proved that for
$\alpha>2$, with arbitrary positive rates, every equilibrium that is not linearly unstable is supported on a forest
\cite[Theorem~3, Corollary~1]{HK}, and that with equal rates and $\alpha>25$ each component of the support is a
whisker \cite[Theorem~2(a)]{HK}, with a threshold independent of the graph. Hirsch, Holmes and Kleptsyn \cite{HHK}
supplied the stochastic input: the vector $N_t/t$ converges almost surely to an equilibrium whose Jacobian has no
positive eigenvalue \cite[Theorem~1]{HHK}, and $E_\infty$ is almost surely the positive support of the limit
\cite[Proposition~1]{HHK}. They also showed that fully supported equilibria on bare paths with at least four edges
are unstable for $\alpha>2$ \cite[Proposition~9]{HHK}, and that the path with three edges carries no positive
equilibrium for $1<\alpha<\alpha_*\approx 4.4$ \cite[Proposition~10(i)]{HHK}.

Kleptsyn \cite[Open Problems~1 and 2, p.~655]{Kl} asked to improve the forest threshold $2$ and the whisker
threshold $25$, and suggested that $4/3$ might be sharp for both; the same two questions are Open Problems~12
and~13 in the collection of Aletti et al.\ \cite{AOP}. Kriebel \cite{Kr} answered the first: for $\alpha>4/3$ and
arbitrary positive rates the set $E_\infty$ is almost surely a forest \cite[Theorem~1(1)]{Kr}, the main step being
the instability of every equilibrium whose support contains an odd cycle of length $m$ when
$\alpha>\sec^2(\pi/(2m))$ \cite[Theorem~3]{Kr}. For the second he proved the whisker conclusion with equal rates
for $\alpha\ge 17/4$ \cite[Theorem~1(2)]{Kr}, leaving the interval $4/3<\alpha<17/4$ open.

Our main theorem closes that interval.

\begin{theorem}\label{thm:main}
Let $G$ be a finite simple loopless graph with equal firing rates and positive integer initial counts, and let
$\alpha>4/3$. Then almost surely every component of $(V,E_\infty)$ with at least one edge is a whisker. If moreover
$\alpha\le 2$, then almost surely every such component is a star.
\end{theorem}

The constant $4/3$ is best possible. For $1<\alpha<4/3$ the uniform equilibrium of the triangle is strictly linearly
stable \cite[Proposition~2(v)]{vdHHKR}, and by \cite[Theorem~1]{HHK} the urn converges to it with positive
probability, so all three edges of the triangle survive with positive probability; this is the observation behind
the suggestion in \cite{Kl,AOP}. At $\alpha=4/3$ the question remains open, since the uniform triangle equilibrium
is then critical. Stars do occur on $(4/3,2]$: the symmetric $k$-edge star, with all weights $(k+1)/k$, is an equilibrium
\cite[Section~3.1]{vdHHKR}, and \Cref{lem:Q} below gives
$Q(u)=\bigl(\frac{\alpha-1}{k}-1\bigr)\sum_e u_e^2-\frac{\alpha}{k^2}\bigl(\sum_e u_e\bigr)^2$, which is negative for
$u\ne 0$ when $\alpha\le 2$ (for $k=1$ it equals $-2u_e^2$), so the Jacobian has only negative eigenvalues.

\Cref{thm:main} follows from the following deterministic statement about equilibria, which we prove in
\Cref{sec:large,sec:small}, together with the stochastic theorems above (\Cref{sec:stoch}).

\begin{theorem}\label{thm:det}
Consider the equal-rate WARM flow on a finite simple graph and an equilibrium whose Jacobian has no positive
eigenvalue.
\begin{enumerate}[label=(\alph*)]
\item If $\alpha>2$, every component of the positive support is a whisker.
\item If $4/3\le\alpha\le 2$, every component of the positive support that is a tree is a star.
\item If $1<\alpha\le 2$, no equilibrium, stable or not, has a support component that is a tree of diameter three.
\end{enumerate}
\end{theorem}

For $\alpha\ge 17/4$, part~(a) is Kriebel's \cite[Sections~4.1--4.3]{Kr}, and we prove it for $2<\alpha\le 6$. Part~(c)
extends \cite[Proposition~10(i)]{HHK} from the bare three-edge path to all trees of diameter three, on the
smaller range $\alpha\le 2$. In part~(b) the support is a forest by Kriebel's theorem \cite[Theorem~3]{Kr}; part~(b)
is where $4/3$ enters a second time, through \Cref{lem:A} below.

The two arguments are different in kind. For $\alpha>2$ the diagonal condition of Holmes and Kleptsyn
\cite[Proposition~1]{HK} says that every edge of the support takes a share at least $1-1/\alpha>1/2$ of the
reinforcement at one of its endpoints. Since no vertex can give more than half to two edges, counting these
endpoints on a tree forces one of two rigid orientations (\Cref{lem:count}). A monotonicity estimate then forbids
long chains (\Cref{lem:depth}), and a polynomial inequality with positive Bernstein coefficients forbids two long
branches (\Cref{lem:branch}). The approach refines the least-weight and champion arguments of \cite{HK}. For
$\alpha\le 2$ the counting breaks down, because $1-1/\alpha\le 1/2$. Here we use the Lyapunov function of
\cite[(A.4)]{vdHHKR}, in the gradient form of \cite[Section~2.2]{HK}, and evaluate its Hessian on a test vector
supported on a path. The increments alternate in sign so that the vertex averages vanish at every inner vertex of the
path. This is the open-chain version of the alternating perturbation around even cycles in
\cite[Lemma~2]{HHK}; on a path, unlike a cycle, two end terms remain, and \Cref{lem:A,lem:B} show that both are
positive when the path runs between two edges next to leaves.

\section{Equilibria and the Hessian form}\label{sec:prelim}

We write $a=(a_e)_{e\in E}\in[0,\infty)^E$ and, for a vertex $v$, $S_v(a)=\sum_{f\ni v}a_f^\alpha$. The equal-rate
WARM flow is
\begin{equation}\label{eq:flow}
  \dot a_e=F_e(a)=-a_e+\sum_{v\in e}\frac{a_e^\alpha}{S_v(a)} ,
\end{equation}
on the set where $S_v>0$ at every non-isolated vertex. At an equilibrium $a$ write $\supp a=\{e: a_e>0\}$ and, for
$e\ni v$,
\[
  q_{e,v}=\frac{a_e^\alpha}{S_v(a)} .
\]
Then $\sum_{e\ni v}q_{e,v}=1$ at every vertex meeting the support, and $a_e=q_{e,x}+q_{e,y}$ for every edge $e=xy$
of the support. We call the equilibrium \emph{stable} if its Jacobian $DF(a)$ has no positive eigenvalue. This
includes critical equilibria, as in \cite[Theorem~1]{HHK}.

For $e\notin\supp a$ and $\alpha>1$ we have $\partial F_e/\partial a_f(a)=-\delta_{ef}$ for every edge $f$, because
$a_e^{\alpha-1}=0$. Hence the eigenvalues of $DF(a)$ are $-1$ together with those of the block $J$ indexed by
$\supp a$, and the coordinates outside the support do not enter $J$. On $(0,\infty)^{\supp a}$ put
\[
  L(a)=-\sum_{e}a_e+\frac1\alpha\sum_v\log S_v(a).
\]
Then $F_e=a_e\,\partial_eL$ \cite[(A.4)]{vdHHKR}, \cite[Lemma~2, Corollary~4]{HK}. At an equilibrium $\partial_eL=0$
on the support, so $J=\diag(a)\,H$ with $H$ the Hessian of $L$. Since $J$ is similar to the symmetric matrix
$\diag(a)^{1/2}H\diag(a)^{1/2}$, its eigenvalues are real, and by Sylvester's law of inertia $J$ has a positive
eigenvalue as soon as $H$ has a positive direction (compare \cite[Corollary~5]{HK}).

\begin{lemma}\label{lem:Q}
At an equilibrium, for $u\in\mathbb R^{\supp a}$ and $w_e=a_eu_e$,
\[
  \langle w,Hw\rangle=Q(u):=\sum_v\Bigl[(\alpha-1)\sum_{e\ni v}q_{e,v}u_e^2-\alpha\Bigl(\sum_{e\ni v}q_{e,v}u_e\Bigr)^2\Bigr],
\]
the sums running over edges of the support. In particular, if $Q(u)>0$ for some $u$, the equilibrium is not stable.
\end{lemma}

\begin{proof}
Differentiating $\partial_eL=-1+\sum_{v\in e}a_e^{\alpha-1}/S_v$ gives
\[
  H_{ef}=\delta_{ef}(\alpha-1)\sum_{v\in e}\frac{q_{e,v}}{a_e^2}
        -\alpha\sum_{v\in e\cap f}\frac{q_{e,v}q_{f,v}}{a_ea_f},
\]
and substituting $w_e=a_eu_e$ gives the formula.
\end{proof}

For an edge $e=xy$ of the support of a stable equilibrium, taking for the test vector the indicator of $e$ gives
$Q\le 0$, that is,
\begin{equation}\label{eq:diag}
  q_{e,x}^2+q_{e,y}^2\ \ge\ \Bigl(1-\frac1\alpha\Bigr)\bigl(q_{e,x}+q_{e,y}\bigr),
\end{equation}
the diagonal condition of Holmes and Kleptsyn \cite[Proposition~1]{HK}; see also \cite[(5)]{Kr}. In particular
$\max(q_{e,x},q_{e,y})\ge 1-1/\alpha$.

The positive support splits into components, and the restriction of a stable equilibrium to one component $C$
satisfies the equilibrium equations on $C$, with $\sum_{e\in C}a_e$ equal to the number of vertices of $C$; a test
vector supported in $C$ sees only $C$. All arguments below are therefore made on one component.

\section{\texorpdfstring{The range $\alpha>2$}{The range alpha > 2}}\label{sec:large}

Throughout this section $\alpha>2$, the equilibrium is stable, and $h=1-1/\alpha\in(1/2,1)$. We say that the
incidence $(e,v)$ is \emph{high} if $q_{e,v}\ge h$. By \eqref{eq:diag} every support edge has a high endpoint, and
since $h>1/2$ and $\sum_{e\ni v}q_{e,v}=1$, no vertex is high for two edges.

The following two facts are due to Holmes and Kleptsyn.

\begin{lemma}[{\cite[Theorem~3, Lemma~4]{HK}; see also \cite[Lemma~4]{Kr}}]\label{lem:hk}
Let $\alpha>2$. Every component of the support is a tree, and every support edge has weight $a_e>1$.
\end{lemma}

\begin{proof}
If a component contains a cycle, let $e=uw$ be an edge of least weight on it. Both $u$ and $w$ carry another cycle
edge of weight at least $a_e$, so $q_{e,u},q_{e,w}\le 1/2<h$, contradicting \eqref{eq:diag}. In a tree component
with at least two edges, the same argument shows that an edge of least weight has a leaf endpoint, where its
incidence is $1$; its other incidence is positive, so its weight exceeds one, and so does every weight in the
component. A component with one edge has weight $2$.
\end{proof}

\begin{lemma}\label{lem:count}
Let $T$ be a component of the support, with $m$ edges. Then exactly one of the following holds.
\begin{enumerate}[label=(\Roman*)]
\item One edge of $T$ is high at both endpoints, and every other edge is high at its endpoint farther from that edge.
\item One vertex $r$ of $T$ is high for no edge, and every edge is high at its endpoint farther from $r$.
\end{enumerate}
\end{lemma}

\begin{proof}
Let $\eta$ be the number of high incidences. Each edge has at least one and each vertex at most one, so
$m\le \eta\le m+1$. If $\eta=m+1$, exactly one edge $e_0$ is high at both endpoints and every vertex is high for exactly one
edge. If $f\ne e_0$ is joined to $e_0$ through a path whose last vertex before $f$ is $x$, then by induction on the
length of the path $x$ is already high for the previous edge, so $f$ is high at its other endpoint. If $\eta=m$, every
edge is high at exactly one endpoint, and exactly one vertex $r$ is high for no edge. Induction on the distance from
$r$ gives the orientation in~(II).
\end{proof}

\begin{lemma}\label{lem:depth}
Let $e=uv$ be a support edge with $q_{e,u}\le 1/\alpha$ and $q_{e,v}\ge h$. Then $v$ is a leaf of the support.
\end{lemma}

\begin{proof}
Suppose $v$ has another support edge $f$. Put $\lambda=q_{e,u}$, $t=1-q_{e,v}\in(0,1/\alpha]$ and $d=q_{f,v}\le t$.
At $v$ we have $d/(1-t)=(a_f/a_e)^\alpha$, and $a_f>1$ by \Cref{lem:hk}, so
$a_e>\bigl((1-t)/t\bigr)^{1/\alpha}$. Since $a_e=\lambda+1-t$,
\[
  \lambda>t-1+\Bigl(\frac{1-t}{t}\Bigr)^{1/\alpha}.
\]
We claim that the right side exceeds $1/\alpha$ for $0<t\le 1/\alpha$, which contradicts $\lambda\le 1/\alpha$. The
claim is $\varphi(t)>0$ for
\[
  \varphi(t)=\log\frac{1-t}{t}-\alpha\log\Bigl(1+\frac1\alpha-t\Bigr).
\]
Differentiating,
\[
  \varphi'(t)=-\frac1{1-t}-\frac1t+\frac{\alpha}{1+1/\alpha-t}
  =\frac{\alpha^2t(1-t)-\alpha-1+\alpha t}{\alpha t(1-t)(1+1/\alpha-t)} .
\]
With $x=\alpha t\in(0,1]$ the numerator equals $\alpha x-x^2-\alpha-1+x=-\bigl(1+(1-x)(\alpha-x)\bigr)<0$, so
$\varphi'<0$ on $(0,1/\alpha]$.
Hence $\varphi(t)\ge\varphi(1/\alpha)=\log(\alpha-1)>0$.
\end{proof}

\begin{lemma}\label{lem:branch}
Let $2<\alpha\le 6$. Let $e=rv$ be a support edge with $q_{e,v}\ge h$, and suppose $v$ has a support edge $f=vy$
with $y$ a leaf. Then $q_{e,r}>1/2$.
\end{lemma}

\begin{proof}
Put $\lambda=q_{e,r}$, $t=1-q_{e,v}\le 1/\alpha$ and $d=q_{f,v}\le t$. Since $y$ is a leaf, $q_{f,y}=1$ and
$a_f=1+d$. At $v$, $d/(1-t)=(a_f/a_e)^\alpha$, so $a_e=(1+d)\bigl((1-t)/d\bigr)^{1/\alpha}$ and
\[
  \lambda=t-1+(1+d)\,d^{-1/\alpha}(1-t)^{1/\alpha}.
\]
The logarithmic derivative of $d\mapsto(1+d)d^{-1/\alpha}$ is $\frac1{1+d}-\frac1{\alpha d}$, which is negative
because $(\alpha-1)d<1$. Hence the right side is smallest at $d=t$, and
\[
  \lambda\ \ge\ t-1+(1+t)\Bigl(\frac{1-t}{t}\Bigr)^{1/\alpha}\ \ge\ t-1+(1+t)\Bigl(\frac{1-t}{t}\Bigr)^{1/6},
\]
where the second step uses $1/\alpha\ge 1/6$ and $(1-t)/t>1$ (as $t\le 1/\alpha<1/2$). The last expression exceeds
$1/2$ if and only if $P(t)=(1-t)(1+t)^6-t(3/2-t)^6>0$. As a polynomial in $\sigma=2t\in[0,1]$, $P$ has degree seven
and Bernstein coefficients
\[
  1,\ \tfrac{487}{896},\ \tfrac{165}{224},\ \tfrac{297}{224},\ \tfrac{243}{112},\ \tfrac{711}{224},\
  \tfrac{1899}{448},\ \tfrac{665}{128},
\]
all positive, so $P>0$ on $[0,1/2]$.
\end{proof}

\begin{proof}[Proof of \Cref{thm:det}(a)]
For $\alpha\ge 17/4$ this is \cite[Sections~4.1--4.3]{Kr}. Let $2<\alpha\le 6$ and let $T$ be a component of the
support; it is a tree by \Cref{lem:hk}. In case~(I) of \Cref{lem:count}, let $e_0=xx'$ be the doubly high edge. Any
other edge $f=xz$ at $x$ has $q_{f,x}\le 1-h=1/\alpha$ and is high at $z$, so $z$ is a leaf by \Cref{lem:depth}; the
same holds at $x'$, and $T$ has diameter at most three. In case~(II), root $T$ at $r$. An edge $vz$ with $v$ at depth
one and $z$ at depth two is high at $z$, and $q_{vz,v}\le 1/\alpha$ because $v$ is high for the edge $rv$; so $z$ is a
leaf by \Cref{lem:depth}, and $T$ has height at most two. If a depth-one vertex $v$ has a child, the edge $e=rv$
satisfies the hypotheses of \Cref{lem:branch}, so $q_{e,r}>1/2$. Since the incidences at $r$ sum to one, at most one
depth-one vertex has children, and $T$ has diameter at most three.
\end{proof}

\begin{remark}
The margin in \Cref{lem:branch} disappears near $\alpha=9$: at $\alpha=9$ and $t=1/12$ the lower bound is below
$1/2$: it equals $-\frac{11}{12}+\frac{13}{12}\,11^{1/9}$, and $11\cdot 13^9<17^9$. This is why the range $\alpha>6$ is taken from \cite{Kr}.
\end{remark}

\section{\texorpdfstring{The range $1<\alpha\le 2$}{The range 1 < alpha <= 2}}\label{sec:small}

In this section the equilibrium is supported on a tree component $T$, and for an edge $e=vx$ of $T$ we put
\[
  d_e(v)=(\alpha-1)\,a_e-\alpha\,q_{e,x}^2 .
\]
The two lemmas below concern an edge $e=cx$ of $T$ such that $c$ has another support edge and $x$ has a leaf
neighbour $y$ in $T$. Write $l=xy$ and
\[
  p=q_{e,x},\qquad s=q_{e,c},\qquad q=q_{l,x}.
\]
Then $q_{l,y}=1$, so $a_l=1+q$; also $a_e=p+s$, $p/q=(a_e/a_l)^\alpha$, $p+q\le 1$, and $s<1$ because $c$ has
another support edge.

\begin{lemma}\label{lem:B}
If $1<\alpha\le 2$, then $p<1/2$.
\end{lemma}

\begin{proof}
Suppose $p\ge 1/2$, so $q\le 1-p$. From $p/q=((p+s)/(1+q))^\alpha$,
\[
  p+s=p^{1/\alpha}(1+q)q^{-1/\alpha}.
\]
The function $q\mapsto(1+q)q^{-1/\alpha}$ is decreasing on $(0,1)$, since its logarithmic derivative
$\frac1{1+q}-\frac1{\alpha q}$ is negative when $(\alpha-1)q<1$. Hence
\[
  p+s\ \ge\ (2-p)\Bigl(\frac p{1-p}\Bigr)^{1/\alpha}\ \ge\ (2-p)\Bigl(\frac p{1-p}\Bigr)^{1/2},
\]
using $p/(1-p)\ge 1$ and $1/\alpha\ge 1/2$. Now
$(2-p)^2p-(1+p)^2(1-p)=(2p-1)(p^2-p+1)\ge 0$, so $p+s\ge 1+p$, that is, $s\ge 1$, a contradiction.
\end{proof}

\begin{lemma}\label{lem:A}
If $\alpha>1$ and $p\le 1/2$, then $s>p$. If moreover $\alpha\ge 4/3$, then $d_e(c)>0$.
\end{lemma}

\begin{proof}
If $s\le p$, then $a_e=p+s\le 2p\le 1<a_l$, so $a_e/a_l<1$ and, since $\alpha>1$,
\[
  \frac pq=\Bigl(\frac{a_e}{a_l}\Bigr)^\alpha<\frac{a_e}{a_l}\le\frac{2p}{1+q},
\]
which gives $q>1$, impossible. Hence $s>p$, and for $\alpha\ge 4/3$
\[
  d_e(c)=(\alpha-1)(p+s)-\alpha p^2>2(\alpha-1)p\ \ge\ \frac\alpha2\,p\ \ge\ \alpha p^2,
\]
where the middle inequality is $4(\alpha-1)\ge\alpha$ and the last one uses $p\le 1/2$.
\end{proof}

The next lemma holds for every $\alpha>1$. The test vector is supported on a path, and its increments alternate so
that the weighted average $\sum_{e\ni v}q_{e,v}u_e$ vanishes at each inner vertex; compare the alternating
perturbation around even cycles in \cite[Lemma~2]{HHK}.

\begin{lemma}\label{lem:path}
Let $y_0y_1\cdots y_{k+1}$ be a path in $T$ with $k\ge 1$ and edges $g_i=y_{i-1}y_i$, $1\le i\le k+1$. Put
$u_{g_1}=1$, $u_{g_{i+1}}=-u_{g_i}\,q_{g_i,y_i}/q_{g_{i+1},y_i}$ for $1\le i\le k$, and $u=0$ on the other edges.
Then
\[
  Q(u)=(\alpha-1)\sum_{i=2}^{k}a_{g_i}u_{g_i}^2+u_{g_1}^2\,d_{g_1}(y_1)+u_{g_{k+1}}^2\,d_{g_{k+1}}(y_k).
\]
\end{lemma}

\begin{proof}
A vertex off the path meets no edge with $u\ne 0$. At an inner vertex $y_i$, $1\le i\le k$, the only such edges are
$g_i$ and $g_{i+1}$, and the choice of $u$ gives $\sum_{e\ni y_i}q_{e,y_i}u_e=0$; so $y_i$ contributes
$(\alpha-1)(q_{g_i,y_i}u_{g_i}^2+q_{g_{i+1},y_i}u_{g_{i+1}}^2)$. The end $y_0$ contributes
$u_{g_1}^2\bigl[(\alpha-1)q_{g_1,y_0}-\alpha q_{g_1,y_0}^2\bigr]$, and $y_{k+1}$ contributes the analogous term.
Collecting by edges and using $a_e=q_{e,x}+q_{e,y}$ for $e=xy$: a middle edge $g_i$, $2\le i\le k$, has both endpoints inner and
collects $(\alpha-1)a_{g_i}u_{g_i}^2$, while $g_1$ collects
$u_{g_1}^2\bigl[(\alpha-1)(q_{g_1,y_1}+q_{g_1,y_0})-\alpha q_{g_1,y_0}^2\bigr]=u_{g_1}^2d_{g_1}(y_1)$, and similarly
for $g_{k+1}$. The vertices of the path are distinct because $T$ is a tree.
\end{proof}

\begin{corollary}\label{cor:facing}
At a stable equilibrium there are no two distinct edges $g=xv$ and $g'=v'x'$ of $T$ with $d_g(v)>0$ and
$d_{g'}(v')>0$ such that the path in $T$ from $x$ to $x'$ begins with $g$ and ends with $g'$ ($v=v'$ is allowed).
\end{corollary}

\begin{proof}
Apply \Cref{lem:path} to the path from $x$ to $x'$: every term is nonnegative, the two end terms are positive, and
\Cref{lem:Q} gives a positive eigenvalue.
\end{proof}

\begin{proof}[Proof of \Cref{thm:det}(c)]
Let $x_1x_2x_3x_4$ be a longest path in a tree component of diameter three. Then $x_1,x_4$ are leaves and the edge
$e=x_2x_3$ has non-leaf endpoints, each with a leaf neighbour. Applying \Cref{lem:B,lem:A} with $c=x_3$, $x=x_2$
gives $q_{e,x_3}>q_{e,x_2}$, and applying them with $c=x_2$, $x=x_3$ gives the reverse inequality. Hence no
equilibrium has such a component. Stability was not used.
\end{proof}

\begin{proof}[Proof of \Cref{thm:det}(b)]
Let $4/3\le\alpha\le 2$ and let $T$ be a tree component of diameter $D$. Diameter three is excluded by part~(c).
Suppose $D\ge 4$ and let $x_1x_2\cdots x_{D+1}$ be a longest path. Its endpoints are leaves, so $x_2,x_3,x_{D-1},x_D$
have degree at least two in $T$. \Cref{lem:B,lem:A} applied to $e=x_3x_2$ (with $c=x_3$, $x=x_2$ and leaf $x_1$)
give $d_e(x_3)>0$, and applied to $e'=x_{D-1}x_D$ (with $c=x_{D-1}$, $x=x_D$ and leaf $x_{D+1}$) give
$d_{e'}(x_{D-1})>0$. As $D\ge 4$ these edges are distinct, and the path $x_2x_3\cdots x_{D-1}x_D$ begins with $e$
and ends with $e'$. \Cref{cor:facing} shows that the equilibrium is not stable. Hence $D\le 2$.
\end{proof}

\begin{example}
Let $\alpha=3/2$ and let $T$ be the path with four edges, with weights $(A,B,B,A)$, where $A=(5+\sqrt5)/4$ and
$B=(5-\sqrt5)/4$. This is a positive equilibrium. The test vector of \Cref{lem:path} on the two middle edges
($k=1$) gives $Q(u)=(7\sqrt5-11)/8>0$, so the equilibrium is unstable, as \Cref{thm:det}(b) requires.
\end{example}

\begin{remark}
\Cref{lem:A} is the only place in this section where a lower bound on $\alpha$ beyond $\alpha>1$ is used, and it
needs only $\alpha\ge 4/3$; \Cref{lem:B} is the only place where $\alpha\le 2$ is used. Hence \Cref{thm:det}(b) holds
for $4/3\le\alpha\le 2$. The stochastic statement at $\alpha=4/3$ fails to follow only because the forest theorem does not
extend there: the uniform triangle equilibrium has eigenvalues $-1,0,0$ at $\alpha=4/3$.
\end{remark}

\section{Proof of the main theorem}\label{sec:stoch}

Start from unit counts. By \cite[Theorem~1]{HHK}, applied to the finite graph $G$ with unit rates, $N_t/t$ converges
almost surely to an equilibrium $a$ of \eqref{eq:flow} whose Jacobian has no positive eigenvalue, and by
\cite[Proposition~1]{HHK}, whose bounded-rate and bounded-degree hypotheses hold on a finite graph,
$E_\infty=\supp a$ almost surely. Kriebel's Lemma~2 \cite{Kr} transfers these almost-sure conclusions to any fixed
positive integer initial counts: a prescribed finite sequence of firings and choices leads from unit counts to the
given counts with positive probability, and the future process then has the law started from those counts.

If $\alpha>2$, every component of $\supp a$ is a whisker by \Cref{thm:det}(a). If $4/3<\alpha\le 2$, the support of
an equilibrium with no positive eigenvalue is a forest by \cite[Theorem~3]{Kr} (an even support cycle is unstable for
every $\alpha>1$, and an odd one of length $m\ge 3$ for $\alpha>\sec^2(\pi/(2m))$, and $\sec^2(\pi/(2m))\le 4/3$),
so every component is a tree and hence a star by \Cref{thm:det}(b). \qed

\begin{remark}
Two questions remain: the case $\alpha=4/3$, and whether $4/3$ is also the whisker threshold for unequal rates.
\Cref{lem:count,lem:depth,lem:branch,lem:A,lem:B} use equal rates through the normalization
$\sum_{e\ni v}q_{e,v}=1$ with a common constant at every vertex and through $a_l=1+q$ at a leaf edge.
\end{remark}

\section*{Acknowledgments and funding}

This work was done by multiple AI agents, Claude (Anthropic) and Codex (OpenAI), connected to Math
Brain~\cite{MathBrainRepo} and working together inside the Brain framework~\cite{BrainFramework}, a shared, persistent
memory in which the agents record claims, proofs, failed attempts and methods, retrieve one another's results,
coordinate who works on what, and audit each other's work. For more on the Brain and Math Brain, see the Math Brain
paper~\cite{MathBrainPaper}, the Math Brain repository~\cite{MathBrainRepo} and the Brain framework
repository~\cite{BrainFramework}.

The research presented in this paper was supported by the European Research Council (ERC) under the European
Union's Horizon Europe research and innovation programme (grant agreement No.\ 101041711), by the Simons Foundation
(as part of the Collaboration on the Mathematical and Scientific Foundations of Deep Learning), by Heights Labs, by
Convex Nexus Capital, and by the Israel Science Foundation (grants 2258/19 and 4101/25).

\section*{Declaration of competing interest}

The author holds equity in Heights Labs and Convex Nexus Capital and has a management role at Convex Nexus Capital.

\section*{Declaration of generative AI and AI-assisted technologies}

The arguments of \Cref{sec:large,sec:small} were developed with AI systems (OpenAI Codex and Anthropic Claude agents
working under the author's direction), which also drafted the text and carried out the exact-arithmetic
computations. They build on the results of van der Hofstad, Holmes, Kuznetsov and Ruszel \cite{vdHHKR}, Holmes and
Kleptsyn \cite{HK}, Hirsch, Holmes and Kleptsyn \cite{HHK} and Kriebel \cite{Kr}, used as cited. The author directed
the work and is responsible for the content of the paper and for decisions concerning its dissemination.

\section*{Data and code availability}

A SymPy script that recomputes the Bernstein coefficients of \Cref{lem:branch}, the factorization in the proof of
\Cref{lem:B}, the derivative in the proof of \Cref{lem:depth} and the example in \Cref{sec:small} is available at
\url{https://github.com/yspennstate/warm-whisker-threshold} (release v1.0).

\appendix
\section{Notation}

\begin{center}
\small
\begin{tabular}{lll}
Symbol & Meaning & First use\\[2pt]
$G=(V,E)$ & finite simple loopless graph & \Cref{sec:intro}\\
$\alpha$ & reinforcement exponent, $\alpha>1$ & \Cref{sec:intro}\\
$N_t(e)$ & count of edge $e$ at time $t$ & \Cref{sec:intro}\\
$E_\infty$ & edges whose counts tend to infinity & \Cref{sec:intro}\\
$a=(a_e)$ & point of $[0,\infty)^E$; an equilibrium & \Cref{sec:prelim}\\
$S_v(a)$ & $\sum_{f\ni v}a_f^\alpha$ & \Cref{sec:prelim}\\
$F$ & vector field of the flow \eqref{eq:flow} & \Cref{sec:prelim}\\
$\supp a$ & positive support $\{e:a_e>0\}$ & \Cref{sec:prelim}\\
$q_{e,v}$ & share $a_e^\alpha/S_v(a)$ of edge $e$ at vertex $v$ & \Cref{sec:prelim}\\
$DF(a)$, $J$ & Jacobian and its block on $\supp a$ & \Cref{sec:prelim}\\
$L$, $H$ & Lyapunov function and its Hessian & \Cref{sec:prelim}\\
$Q(u)$ & Hessian form in the variables $u_e=w_e/a_e$ & \Cref{lem:Q}\\
$h$ & $1-1/\alpha$, threshold for a high incidence & \Cref{sec:large}\\
$\eta$ & number of high incidences of a tree component & \Cref{lem:count}\\
$T$ & a tree component of the support & \Cref{lem:count}\\
$d_e(v)$ & $(\alpha-1)a_e-\alpha q_{e,x}^2$ for $e=vx$ & \Cref{sec:small}\\
$p$, $s$, $q$ & $q_{e,x}$, $q_{e,c}$, $q_{l,x}$ for $e=cx$ and leaf edge $l=xy$ & \Cref{sec:small}\\
\end{tabular}
\end{center}

\begin{thebibliography}{99}

\bibitem{AOP}
G.~Aletti, A.~Collevecchio, I.~Crimaldi, M.~Deijfen, S.~Dereich, A.~Ghiglietti, R.~van der Hofstad, M.~Holmes,
D.~Kious, V.~Kleptsyn, B.~Lodewijks, P.~M\"orters and B.~Schapira,
\emph{Open problems in stochastic reinforcement models},
MATRIX Annals, 2023.
\url{https://www.matrix-inst.org.au/wp_Matrix2016/wp-content/uploads/2023/Aletti.pdf}

\bibitem{HHK}
C.~Hirsch, M.~Holmes and V.~Kleptsyn,
\emph{Infinite WARM graphs III: strong reinforcement regime},
Nonlinearity \textbf{36} (2023), 3013--3042.
\doi{10.1088/1361-6544/acc9a0}

\bibitem{HK}
M.~Holmes and V.~Kleptsyn,
\emph{Proof of the WARM whisker conjecture for neuronal connections},
Chaos \textbf{27} (2017), 043104.
\doi{10.1063/1.4978683}

\bibitem{Kl}
V.~Kleptsyn,
\emph{Graph-based interacting P\'olya urns},
Oberwolfach Reports \textbf{20} (2023), no.~12, 653--655.
\doi{10.4171/OWR/2023/12}

\bibitem{Kr}
A.~Kriebel,
\emph{A sharp forest bound and an improved whisker bound for finite WARM graphs},
preprint, version 1.0.0, 2026.
\doi{10.5281/zenodo.23592692}

\bibitem{MathBrainRepo}
Y.~Shmalo, \emph{Math Brain: public snapshot, code and tests}, 2026.
\url{https://github.com/yspennstate/math-brain-public}

\bibitem{MathBrainPaper}
Y.~Shmalo, \emph{Math Brain: shared, auditable memory for multi-agent mathematical research}, October 2026.
\url{https://github.com/yspennstate/math-brain-public/tree/main/papers/math-brain-2026-10-24}

\bibitem{BrainFramework}
Y.~Shmalo, \emph{Brain Framework}, public software repository, 2026.
\url{https://github.com/yspennstate/brain-framework-public}

\bibitem{vdHHKR}
R.~van der Hofstad, M.~Holmes, A.~Kuznetsov and W.~Ruszel,
\emph{Strongly reinforced P\'olya urns with graph-based competition},
Ann.\ Appl.\ Probab.\ \textbf{26} (2016), 2494--2539.
\doi{10.1214/16-AAP1153}; arXiv:1406.0449.

\end{thebibliography}

\end{document}
